%% kapa: exam=1 round=32 year=2021 period=1 form=A subject=civil_law subject_ko=민법 seq=4 number=4
%% source: Q-Net 감정평가사 기출문제 (cst00309) artlSeq=5208397 file=extracted/1008/5208397/A형/2021년도 제32회 감정평가사 1차 1교시 A형.pdf
%% conversion: kordoc 4.24.51 -> tools/gen1.py (auto); review=proofread-sonnet; flags=none
\begin{question}{4}
실종선고에 관한 설명으로 옳지 않은 것은? (다툼이 있으면 판례에 따름)
\begin{choices}
  \item 가족관계등록부상 이미 사망으로 기재되어 있는 자에 대해서는 원칙적으로 실종선고를 할 수 없다.
  \item 실종선고를 받아 사망으로 간주된 자는 실종선고가 취소되지 않는 한 반증을 통해 그 효력을 번복할 수 없다.
  \item 실종선고 후 그 취소 전에 선의로 한 행위의 효력은 실종선고의 취소에 의해 영향을 받지 않는다.
  \item 실종선고의 취소에는 공시최고를 요하지 않는다.
  \item 실종자를 당사자로 한 판결이 확정된 후에 실종선고가 확정되어 그 사망간주의 시점이 소 제기 전으로 소급하는 경우, 특별한 사정이 없는 한 그 판결은 당사자능력이 없는 사람을 상대로 한 판결로서 무효가 된다.
\end{choices}
\end{question}
\answer{5}
